\chapter*{مقدمه}
\addcontentsline{toc}{chapter}{مقدمه}
\label{ch:intro}

این گزارش ادامه گزارش مرحله اول طرح است. در مرحله اول فاز مطالعاتی و نتایج اولیه کار ارائه شد. نخست، برای آشنایی با ابزارها، نتایج یک مطالعه تحلیل بوزون هیگز با جداکننده کوانتومی بازتولید شد: داده‌ها با یک خودرمزگذار کلاسیک از ۶۷ ویژگی به فضای نهان ۸ یا ۱۶ بعدی فشرده و با ماشین بردار پشتیبان کوانتومی (\lr{QSVM}) و مدار کوانتومی وردشی با بارگذاری مجدد داده طبقه‌بندی شدند، و مدل‌های کوانتومی عملکردی هم‌سطح مدل‌های کلاسیک مشابه داشتند. سپس شناساگرهای جت با طعم سنگین و جت‌های خیزبافته و الگوریتم‌های کوانتومی موجود برای آنها مرور شد. بر این پایه، برای طراحی شناساگر کوانتومی جت‌های بوست‌شده یک نمونه ۲۰۰۰۰ جتی پیش‌پردازش شد، روش‌های کاهش بعد (\lr{PCA} و سه نوع خودرمزگذار) مقایسه شدند و یک شبکه عصبی گرافی کوانتومی روی فضای فشرده‌شده، با آموزش جداگانه، ارزیابی شد. در کنار این کارها زیرساخت محاسباتی، یعنی نرم‌افزارهای شبیه‌سازی و خوشه محاسباتی \lr{CERN}، فراهم شد. در پایان آن گزارش دو گام بعدی مشخص شده بود: آموزش هم‌زمان خودرمزگذار و مدار کوانتومی، و تولید داده برای جداسازی رویدادهای \ttZ از پس‌زمینه \ttbb.

هدف این گزارش انجام همین دو گام است. در فصل~\ref{ch:tagger} رمزگذار کلاسیک و شبکه عصبی گرافی کوانتومی هم‌زمان آموزش داده می‌شوند تا جت‌های \Zbos{} بوست‌شده از جت‌های گلوئونی جدا شوند، و عملکرد آن با \lr{ParticleNet} مقایسه می‌شود (بخشی از این کار در مقاله‌ای گزارش شده که در آخرین مراحل داوری نشریه \lr{EPJC} است)؛ بخش~\ref{sec:ch1-link} همان فصل پیوند این کار را با گزارش مرحله اول نشان می‌دهد. در فصل~\ref{ch:data} داده مناسب برای جداسازی رویدادهای \ttZ از \ttbb فراهم می‌شود: همه رویدادها، یعنی نمونه‌های سیگنال \Zjet و \ttZ و نمونه‌های پس‌زمینه \bbbb و \ttbb، در یک چارچوب شبیه‌سازی واحد تولید شده‌اند تا تفاوت توزیع‌ها فقط از فیزیک فرایندها ناشی شود و نه از روش تولید. سپس متغیرهای جت و ذره برای هر چهار نمونه با هم مقایسه می‌شوند و بررسی می‌شود که آیا شناساگری که روی یک جفت (سیگنال، پس‌زمینه) آموزش دیده روی جفت دیگر هم کار می‌کند.
